\documentclass[a4paper, 11pt]{article} % das Papierformat zuerst
\usepackage[utf8]{inputenc}
\usepackage{geometry}
\geometry{a4paper, top=25mm, left=35mm, right=25mm, bottom=25mm}
\usepackage{graphicx}
\usepackage{color}
\usepackage{epstopdf}
\usepackage[T1]{fontenc}
\usepackage{setspace}
\usepackage{textcomp}
\usepackage{tabularx}
\usepackage{blindtext}
\usepackage[ngerman]{babel} %deutsche Silbentrennung
\usepackage{titlesec}
\usepackage{enumitem}
\usepackage{ifthen}
\usepackage[colorlinks=true,linkcolor=fsi]{hyperref}

\definecolor{fsi}{RGB}{0,81,150}

\newcounter{summe}
\newcommand{\abstimmung}[4]{
	\setcounter{summe}{#3}
	\addtocounter{summe}{#4}
	\begin{flushleft}
		#1\\
	Es wird über den Antrag abgestimmt.
	\end{flushleft}
	\ifthenelse{\equal{#3}{0}\AND\equal{#4}{0}}{
	\begin{center}
		\textbf{Der Antrag wird einstimmig angenommen.}
	\end{center}
	}{
	\begin{center}
		#2 \ifthenelse{\equal{#2}{1}}{Stimme}{Stimmen} dafür, #3 \ifthenelse{\equal{#3}{1}}{Stimme}{Stimmen} dagegen, #4 \ifthenelse{\equal{#42}{1}}{Enthaltung}{Enthaltungen}\\
		\ifthenelse{#2>\value{summe}}{
		\textbf{Der Antrag ist somit angenommen.}
		}{
		\textbf{Der Antrag ist somit abgelehnt.}
		}
	\end{center}
	}
}


% VARIABLEN
\newcommand{\protokoller}{Philipp Dippel}
\newcommand{\dateOfMeeting}{29.07.2015}
\newcommand{\TeXer}{Philipp Dippel}
\newcommand{\fsiPresident}{Georg Schäfer}
\newcommand{\protocolID}{9}

\begin{document}
%deckblatt start

\doublespacing
\thispagestyle{empty}

\begin{center}
\includegraphics[width=10.0cm]{../logo_fsi.eps}

\vspace*{\fill}
{\LARGE \textbf{Protokoll der Sitzung des Fachschaftsrates \\vom \dateOfMeeting}}\\
Fachschaftsrat Informatik\\
Trier University of Applied Sciences\\
\vspace{2.5cm}
\textit{
	Protokollführer: \textbf{\protokoller} \\
	\LaTeX - Umsetzung von \TeXer\\
	am \today\\
}
Protokoll ID: \textbf{\#\protocolID}
\vfill
\end{center}

\hspace*{-35cm}
\textcolor{fsi}{\rule{64.9cm}{15pt}}
\pagebreak

\setcounter{tocdepth}{2}
\tableofcontents
\pagebreak

\section{Eröffnung}
Als Protokollführer wird \protokoller~bestimmt.\\
Der Sprecher des Fachschaftsrates \fsiPresident~eröffnet die Sitzung um 19:17 Uhr.
\\\\
\textbf{Es wird festgestellt, dass der Fachschaftsrat beschlussfähig ist.}
\textbf{Es fehlen:} Tobias Meier, Fabio Gimmillaro und Michael Ochmann.
Marco Lochen tritt der Sitzung um 19:21 bei.
\\\\
Das Protokoll der letzten Sitzung wurde verlesen und genehmigt.

\section{Events}
Die Eventmanager verkünden, dass es keine nennenswerten Neuigkeiten von vergangenen Events gibt und keine weiteren Events in näherer Zeit anstehen.

\section{Finanzen}
Der Finanzreferent Tim Grundmanns und sein Stellvertreter Philipp Dippel verkünden die Endgeldbestände der Netzwerkinstallationsveranstaltung (NIV) vom 26. Juli 2015. \\ Es ist mit bedauern zu Verkünden, dass nach Abzug der Getränkekosten vom Lieferanten Herrestahl, der Fachschaftsrat Informatik einen Verlust von ca. 100,00 Euro vorzuweisen hat. Aufgrund der Verkaufszahlen müsste ein Kassenbestand mit einem Plus von ca. 100,00 Euro vorherschen. Die negative Bilanz der Kasse lässt sich nach Aussage der Finanzreferenten nur dadurch erklären, dass es bereits während der NIV zu einem unerlaubtem Zugriff auf die Kasse kam. Woher dieser Zugriff kam, lässt sich im Nachhinein nicht mehr Nachweisen. Ein Zugriff auf die  Kasse nach der Veranstaltung kann, aufgrund der Tatsache, dass die Kasse für die Öffentlichkeit unzugänglich gemacht und bewacht wurde, ausgeschlossen werden.\\\\ Die Fachschaft ist sich daher einig, dass auf Zukünftigen Veranstaltung mit Verkauf verschärfte Kontrollen beim Umgang mit dem Kassengeschäft von Nöten sind.
\\\\
Weitergehend verkündet der Finanzreferent Tim Grundmanns, dass dieser in das nächste Studierendenparlament (StuPa) gewählt wurde und somit laut Satzung der Fachschaft Informatik, bis zur konstituierenden Sitzung des StuPa, von seinem Amt zurücktreten und aus dem Fachschaftsrat austreten muss.
\\\\ \textbf{Tim Grundmanns tritt von seinem Amt als Finanzreferent zurück. Ebenso tritt auch Philipp Dippel von seinem Amt der stellvertretender Finanzreferent zurück.}
\\

\noindent Aufgrund dessen müssen diese beiden Ämter neu besetzt werden. Die Wahl findet öffentlich durch Aklamation statt.
\vspace{1em}
\par\noindent
Zur Wahl für das Amt des Finanzreferenten stellt sich Philipp Dippel.
\begin{center}
	\textbf{Die Wahl wird einstimmig angenommen.}
\end{center}

\par\noindent
Zur Wahl für das Amt des stellvertretenden Finanzreferenten stellt sich Elias Broschin.
\begin{center}
	\textbf{Die Wahl wird einstimmig angenommen.}
\end{center}

\par\noindent
Beide Kandidaten nehmen die Wahl an.

\section{Raum}
Der Raumwart Benjamnin Albsmeier verkündet, dass mittlerweile die Termine zur Reinigung des Fachschaftsraumes besser als zuvor eingehalten werden, es aber noch Potential zur Verbesserung gibt.

\section{Server}
Der Administrator des Fachschaftsrates, Johannes Kirchner, gibt bekannt, dass die Inbetriebnahme des neuen Servers der Fachschaft ohne Probleme von statten ging. Einige der geplanten Dienste laufen bereits. Darunter zum Beispiel: GitLab, Minecraft (beide noch nicht öffentlich zugänglich) und OwnCloud.

\section{Sonstiges}
Der Sprecher Georg Schäfer verkündet, dass er sich beim Algemeinen Studierenden Ausschuss (AStA) bezüglich einer Aufwandsentschädigung für die Helfer, die am Entwurf der Satzungsänderung beteiligt waren, erkundigt und diese bewilligt bekommen hat.

\abstimmung{Der Sprecher stellt den Antrag zur Festlegung der Höhe der Aufwandsentschädigung, für die Helfer des Satzungsänderungsausschuss, auf 60 Euro.}{10}{0}{0} %-- {dafür}{dagegen}{enthaltungen}

\par\noindent
Es ergeht der Beschluss (\#FSRI\_2) die Höhe der Aufandsentschädigung auf 60 Euro festzulegen.

\vspace{1em}
\par\noindent
Außerdem erwähnt der Sprecher, dass die Neuwahlen des Fachschaftsrats, welche im kommenden Semester anstehen, nun nicht mehr auf der Vollversammlung abgehalten werden müssen, sondern nach der Wahlordnung der Studierendenschaft durchgeführt werden können. Dadurch soll eine übersichtlichere Durchführung der Wahlen möglich sein.
\\\\
\noindent \textbf{Marco Lochen tritt von seinem Amt als Kommunikationsminister zurück. Er tritt weiter aus dem Fachschaftsrat Informatik der Hochschule Trier aus.}
\\\\
\noindent Der Sprecher Georg Schäfer schließt die Sitzung um 19:45 Uhr.
\\\\\\
\pagebreak
\section{Bestätigung des Protokolls}
Der Sprecher des Fachschaftsrates Informatik sowie der Protokollführer dieses Protokolls bestätigen mit Ihrer Unterschrift unter diesem Protokoll, dass selbiges inhaltlich korrekt ist und alle darin aufgeführten Beschlüsse so wie beschrieben vom Rat getragen werden.
\\

\vspace{3.5cm}
\hrulefill \hfill \hrulefill

\fsiPresident \hfill \protokoller

{\footnotesize (Sprecher)\hfill (Protokollführer)}
\end{document}
